% !TEX root = ../main.tex
\chapter{Tổng quan}
\label{chap:intro}

% ============================================================
\section{Bối cảnh và sự cần thiết của đề tài}
\label{sec:background}

Trong lịch sử phát triển của lĩnh vực Xử lý Ngôn ngữ Tự nhiên (Natural Language
Processing -- NLP), các mô hình dạng Encoder-only như BERT~\cite{devlin2019bert} và
RoBERTa~\cite{liu2019roberta} từng đóng vai trò nền tảng không thể thiếu cho các bài
toán hiểu ngôn ngữ (Natural Language Understanding -- NLU) như phân loại văn bản,
nhận dạng thực thể có tên, suy luận ngôn ngữ tự nhiên và trả lời câu hỏi. Nhờ cơ
chế mã hóa hai chiều (bidirectional encoding), các mô hình này có khả năng nắm bắt
ngữ cảnh từ cả hai phía của một chuỗi đầu vào, tạo ra biểu diễn ngữ nghĩa phong phú
và được ứng dụng rộng rãi trong thực tế.

Tuy nhiên, trong khi các dòng mô hình Decoder-only -- hay còn gọi là các Mô hình
Ngôn ngữ Lớn (Large Language Models -- LLMs) như LLaMA~\cite{dubey2024llama} và
DeepSeek~\cite{deepseek2025} -- liên tục tiến hóa với quy mô dữ liệu và những đột
phá về kiến trúc, các bộ Encoder-only lại gần như không có sự cải thiện đáng kể trong
suốt nhiều năm qua. BERT và RoBERTa vẫn tiếp tục được sử dụng rộng rãi dù kiến thức
của chúng ngày càng trở nên lỗi thời, đặc biệt trong các ứng dụng đòi hỏi hiểu ngữ
cảnh dài hoặc tốc độ suy luận cao.

Sự ra đời của NeoBERT~\cite{lebreton2025neobert} vào năm 2025 đã thu hẹp đáng kể
khoảng cách này. NeoBERT tái định nghĩa khả năng của các mô hình Encoder bằng cách
tích hợp các cải tiến kiến trúc tiên tiến nhất từ thế giới LLM: hàm kích hoạt
(activation function) SwiGLU, cơ chế mã hóa vị trí tương đối Rotary Position Embeddings (RoPE), chuẩn hóa
Pre-RMSNorm, và tỷ lệ chiều sâu trên chiều rộng tối ưu -- sâu 28 lớp nhưng chỉ
giữ kích thước ẩn (hidden size) ở mức 768. Với khoảng 250 triệu tham số và được huấn luyện trên 2,1 nghìn tỷ token,
NeoBERT vượt trội BERT\textsubscript{large} và RoBERTa\textsubscript{large} trên chuẩn
đánh giá MTEB (Massive Text Embedding Benchmark), đồng thời hỗ trợ ngữ cảnh lên
tới 4.096 token -- gấp 8 lần so với BERT và RoBERTa (512 token).

Tuy nhiên, NeoBERT chỉ được huấn luyện trên ngữ liệu tiếng Anh. Đối với tiếng Việt,
nhu cầu về một bộ Encoder mạnh mẽ và hiện đại ngày càng cấp thiết khi các ứng dụng
NLP trong nước -- từ hệ thống tìm kiếm ngữ nghĩa, phân tích cảm xúc, đến các hệ
thống hỏi-đáp tự động -- đòi hỏi chất lượng biểu diễn ngữ nghĩa ngày càng cao. Các
mô hình Encoder tiếng Việt hiện có như PhoBERT~\cite{nguyen2020phobert} và
ViDeBERTa~\cite{lai2023videberta} tuy hiệu quả nhưng vẫn dựa trên kiến trúc
Transformer cũ, bị giới hạn ở độ dài ngữ cảnh 512 token và thiếu đi các tối ưu hóa
kiến trúc hiện đại.

Việc huấn luyện lại NeoBERT từ đầu (from scratch) cho tiếng Việt trên quy mô hàng
nghìn tỷ token là điều bất khả thi về mặt tài nguyên tính toán đối với hầu hết các
nhóm nghiên cứu. Bên cạnh đó, phương pháp tiếp tục tiền huấn luyện (Continued
Pre-training) thông thường -- tức là khởi tạo ngẫu nhiên lớp véc-tơ mới rồi huấn
luyện trực tiếp -- thường dẫn đến hiện tượng \textit{mất thông tin} (catastrophic
forgetting), làm phá vỡ các đặc trưng biểu diễn quý giá mà mô hình đã học từ dữ
liệu tiếng Anh.

Xuất phát từ những thách thức trên, đề tài nghiên cứu phương pháp chuyển đổi ngôn
ngữ hiệu quả cho NeoBERT thông qua kỹ thuật Chuyển đổi tuyến tính trong không gian
có ngữ nghĩa (Semantic Aware Linear Transfer -- SALT)~\cite{lee2025salt}. Thay vì huấn luyện từ đầu, SALT tái sử dụng
(recycling) bộ véc-tơ từ mô hình tiền huấn luyện tiếng Việt chất lượng cao
ViDeBERTa, chuyển đổi không gian biểu diễn sang không gian của NeoBERT thông qua
các phép ánh xạ tuyến tính cục bộ dựa trên độ tương đồng ngữ nghĩa. Đây là một hướng tiếp cận
tiết kiệm tài nguyên, vừa kế thừa không gian biểu diễn tiếng Việt sẵn có, vừa
bảo toàn được sức mạnh kiến trúc của NeoBERT.

% ============================================================
\section{Mục tiêu nghiên cứu}
\label{sec:objectives}

Đề tài đặt ra các mục tiêu nghiên cứu cụ thể như sau:

\begin{enumerate}
    \item \textbf{Nghiên cứu và nắm vững kiến trúc NeoBERT}: Phân tích chi tiết các
    cải tiến kiến trúc của NeoBERT so với BERT/RoBERTa truyền thống, bao gồm SwiGLU,
    RoPE, Pre-RMSNorm và tỷ lệ chiều sâu trên chiều rộng tối ưu; từ đó xác định những điểm cần
    điều chỉnh khi thích ứng sang tiếng Việt.

    \item \textbf{Áp dụng và cải tiến kỹ thuật SALT để chuyển giao không gian biểu
    diễn}: Nghiên cứu và triển khai thuật toán SALT -- bao gồm trích xuất tập điểm neo
    (anchor extraction) dựa trên độ tương đồng ngữ nghĩa và ánh xạ tuyến tính cục bộ
    theo phương pháp bình phương tối thiểu -- nhằm ánh xạ không gian véc-tơ từ
    ViDeBERTa sang không gian véc-tơ của NeoBERT. Trên nền tảng đó, đề tài đề xuất ba
    cải tiến so với SALT gốc: mở rộng tập điểm neo bằng các cặp dịch Anh--Việt, khởi 
    tạo riêng lớp đầu ra kèm hệ số điều chỉnh theo tần suất token,
    và bổ sung giai đoạn căn chỉnh với phần thân đóng băng.

    \item \textbf{Xây dựng mô hình NeoBERT cho tiếng Việt (ViNeoBERT)}: Kết hợp lớp véc-tơ đầu vào và lớp đầu ra đã được khởi tạo bằng SALT với phần thân NeoBERT gốc, sau đó huấn
    luyện qua giai đoạn căn chỉnh với phần thân đóng băng và tiếp tục tiền huấn luyện
    (Continued Pre-training) trên ngữ liệu tiếng Việt quy mô lớn theo lịch tốc độ
    học Warmup--Stable--Decay.

    \item \textbf{Đánh giá hiệu quả của phương pháp}: Đánh giá ViNeoBERT trên một bộ
    tác vụ hiểu ngôn ngữ tiếng Việt đa dạng -- gồm hỏi đáp đọc hiểu (UIT-ViQuAD), suy
    luận, phân loại văn bản, gán nhãn chuỗi, tương đồng ngữ nghĩa và biểu diễn câu
    (Retrieval, Reranking) -- và đối chiếu với các mô hình tiếng Việt và đa ngôn ngữ
    hiện có (PhoBERT, XLM-R). Bên cạnh đó, đề tài so sánh các phương án khởi tạo khác nhau
    khi cùng được huấn luyện trên một lượng token như nhau, nhằm đo riêng phần đóng góp
    của từng thành phần trong phương pháp.

    \item \textbf{Kiểm chứng tính hiệu quả của việc tái sử dụng không gian biểu diễn}: Thông qua
    kết quả thực nghiệm, xác minh rằng việc khởi tạo véc-tơ bằng kỹ thuật SALT -- vốn
    dựa trên không gian biểu diễn sẵn có của ViDeBERTa -- kết hợp với giai đoạn căn chỉnh
    khi phần thân bị đóng băng giúp giảm đáng kể lượng token cần huấn luyện và giữ cho
    quá trình thích ứng ổn định, thay vì phải học lại gần như từ đầu như khi khởi tạo
    ngẫu nhiên.
\end{enumerate}

% ============================================================
\section{Phạm vi nghiên cứu}
\label{sec:scope}

Đề tài tập trung nghiên cứu thực nghiệm trong các giới hạn sau:

\textbf{Đối tượng nghiên cứu chính} bao gồm: (1)~kiến trúc NeoBERT với trọng số
gốc được công bố công khai; (2)~thuật toán SALT với cơ chế trích xuất tập điểm neo và
ánh xạ tuyến tính cục bộ; (3)~mô hình nguồn véc-tơ (donor) ViDeBERTa cung cấp không
gian véc-tơ tiếng Việt cho bước chuyển giao.

\textbf{Dữ liệu thực nghiệm} gồm phần tiếng Việt của kho ngữ liệu web đa ngôn ngữ
CulturaX cho giai đoạn tiếp tục tiền huấn luyện, cùng một bộ tác vụ đánh giá tiếng Việt
đa dạng cho bước tinh chỉnh cho tác vụ ứng dụng -- từ hỏi đáp đọc hiểu (UIT-ViQuAD), suy luận
ngôn ngữ tự nhiên (XNLI, ViANLI, QNLI), phân loại văn bản (UIT-VSFC, UIT-VSMEC,
UIT-ViCTSD), đánh giá tính hợp ngữ pháp (ViCoLA-syn), gán nhãn từ loại (UD-VTB),
tương đồng ngữ nghĩa (STS, Paraphrase), tới biểu diễn câu (Retrieval, Reranking).

\textbf{Các ràng buộc và giới hạn}: Nghiên cứu \textit{không} thực hiện huấn luyện
NeoBERT từ đầu do hạn chế tài nguyên tính toán. Phạm vi bài toán chỉ bao gồm các
tác vụ NLU và trích xuất đặc trưng (Feature Extraction), không bao gồm các bài toán
sinh văn bản tự do của các mô hình Decoder-only. Mô hình kế thừa kiến trúc hỗ trợ
ngữ cảnh tới 4.096 token của NeoBERT và được huấn luyện bổ sung một pha ngắn ở độ
dài này; việc kiểm chứng ngữ cảnh dài chỉ dừng ở một độ đo nội tại -- tức đo trực tiếp
trên bản thân mô hình ngôn ngữ mà không thông qua tác vụ ứng dụng, cụ thể là
pseudo-perplexity theo độ dài chuỗi -- còn đánh giá chuyên biệt cho ngữ cảnh dài
được để lại như hướng phát triển.

% ============================================================
\section{Đóng góp của đề tài}
\label{sec:contributions}

Đề tài đưa ra các đóng góp cụ thể sau:

\begin{itemize}
    \item \textbf{Mô hình NeoBERT cho tiếng Việt (ViNeoBERT)}: Cung cấp một bộ Encoder
    hiện đại khoảng 250 triệu tham số cho tiếng Việt, kế thừa các cải tiến kiến trúc
    của NeoBERT (RoPE, SwiGLU, Pre-RMSNorm, FlashAttention) cùng trần ngữ cảnh
    4.096 token, tương thích với thư viện Hugging Face Transformers và sẵn sàng sử
    dụng cho các bài toán NLU thực tế.

    \item \textbf{Quy trình chuyển đổi ngôn ngữ}: Đề xuất và kiểm chứng một
    quy trình đưa NeoBERT sang tiếng Việt -- gồm khởi tạo véc-tơ, căn chỉnh với phần thân
    đóng băng và tiếp tục tiền huấn luyện -- lần đầu tiên áp dụng cho một kiến trúc Encoder
    thế hệ mới. Ở giai đoạn khởi tạo, đề tài kết hợp kỹ thuật SALT với ba cải tiến so với
    nghiên cứu gốc: (1)~mở rộng tập điểm neo bằng các cặp dịch Anh--Việt khai thác tự động;
    (2)~khởi tạo riêng lớp đầu ra, kèm hệ số điều chỉnh theo tần suất token và hệ số
    tỉ lệ trọng số giúp phân phối tần suất chi phối giai đoạn huấn luyện đầu; (3)~bổ sung giai
    đoạn căn chỉnh với phần thân đóng băng trước khi tiếp tục tiền huấn luyện.

    \item \textbf{Bộ mã nguồn công khai}: Bao gồm thuật toán trích xuất tập điểm neo,
    phần cài đặt kỹ thuật SALT cho lớp véc-tơ đầu vào và lớp đầu ra, cùng kịch bản huấn luyện
    thích ứng, phục vụ cho cộng đồng nghiên cứu NLP tiếng Việt.

    \item \textbf{Đóng góp khoa học}: Kết quả thực nghiệm khẳng định tính khả thi
    của việc tái sử dụng không gian biểu diễn từ các mô hình sẵn có để chuyển một
    kiến trúc Encoder thế hệ mới sang một ngôn ngữ khác với lượng token huấn luyện
    nhỏ hơn nhiều so với huấn luyện lại từ đầu. Hướng
    tiếp cận này tiết kiệm chi phí cho việc cập nhật, áp dụng các kiến trúc thế hệ
    mới, và đặc biệt có ý nghĩa với những ngôn ngữ ít tài nguyên.
\end{itemize}

% ============================================================
\section{Cấu trúc đề tài}
\label{sec:structure}

Phần còn lại của khóa luận được tổ chức như sau:

\begin{itemize}
    \item \textbf{Chương 1 -- Tổng quan}: Giới thiệu bối cảnh và sự cần thiết của
    đề tài, mục tiêu nghiên cứu, phạm vi nghiên cứu và các đóng góp chính.

    \item \textbf{Chương 2 -- Cơ sở lý thuyết và các công trình liên quan}: Trình
    bày kiến trúc Transformer và dòng mô hình Encoder-only, các đặc điểm chính của
    NeoBERT, cơ sở hình học của không gian biểu diễn, cùng các phương pháp khởi tạo
    véc-tơ khi chuyển một mô hình sang ngôn ngữ mới, dẫn tới kỹ thuật SALT.

    \item \textbf{Chương 3 -- Phương pháp xây dựng NeoBERT cho tiếng Việt}: Trình bày
    chi tiết quy trình xây dựng của đề tài: giảm kích thước từ vựng và trích xuất tập
    điểm neo, kỹ thuật SALT cho cả lớp véc-tơ đầu vào lẫn lớp đầu ra, giai đoạn căn chỉnh với phần thân đóng băng và tiếp tục tiền huấn luyện theo lịch WSD.

    \item \textbf{Chương 4 -- Thực nghiệm và Đánh giá}: Mô tả dữ liệu, cấu hình
    huấn luyện, các độ đo và phân tích kết quả thực nghiệm, bao gồm so sánh các
    phương án khởi tạo và đánh giá trên tác vụ ứng dụng.

    \item \textbf{Chương 5 -- Kết luận và Hướng phát triển}: Tổng kết các kết quả
    đạt được so với mục tiêu đề ra và đề xuất các hướng phát triển tiếp theo.
\end{itemize}
